\subsection{External Paid-Parking Validation}
\label{subsec:parking_validation}

\textit{Design.}
We evaluate whether POI forecasts improve prediction and monitoring of payment-based parking pressure. SDOT reports paid occupancy inferred from transactions \citep{sdot_parking_data}; we divide this count by recorded parking capacity. For blockface $\ell$ and clock time $\nu\in\{10{:}00,14{:}00,18{:}00\}$, let $\omega_{b,h,\ell,\nu}$ denote this ratio at the target date of sample $b$ and lead $h$. Ratios above one remain in the response. The eligible blockface--clock-time pairs are fixed from training observations, with 1,343, 1,341, and 885 candidates at the three clock times.

The parking-only ridge model uses a smoothed training mean for each blockface, weekday, and clock time, together with seven- and fourteen-day occupancy lags, lag-missingness indicators, weekday, and clock time. Smoothing adds five observations at the training blockface and clock-time mean. Let $\mathbf u_{b,h,\ell,\nu}$ collect these standardized predictors. The augmented model adds predicted visits within a catchment:
\begin{equation}
\widehat\omega_{b,h,\ell,\nu}
=\max\left\{0,\beta_0+\boldsymbol\beta^{\mathsf T}\mathbf u_{b,h,\ell,\nu}
+\beta_{\mathrm{POI}}\log\left(1+\sum_{j\in\mathcal N_R(\ell)}\hat Y_{b,h,j}\right)\right\},
\label{eq:parking_mapping}
\end{equation}
where $\mathcal N_R(\ell)$ contains POIs within great-circle distance $R$ of blockface $\ell$, and the $\beta$ coefficients are fitted transport-mapping parameters. The added term represents predicted activity around the blockface. The primary radius is 500 m, with 250 and 1,000 m as sensitivities; all blockfaces have matched POIs at every radius.

The mapping uses each observed training date and clock time once, through November 19. Training substitutes contemporaneous observed visits for $\hat Y$ in Eq.~\eqref{eq:parking_mapping}; validation and test substitute saved POI forecasts. The ridge penalty is selected from $\{10^{-4},10^{-2},1,100\}$ by validation occupancy-ratio MAE at lead seven. January retains all coefficients and selections. The primary lead is seven days, with one- and three-day sensitivities. The POI inputs use $T=14$, MAE training, and all five seeds. Calendar and seven-day-lag parking forecasts provide additional comparators.

A monitoring budget selects the highest-scored 10\% or 20\% of the training-defined candidates at each target date and clock time. Selection precedes the outcome-availability mask; observed selected locations determine precision, and observed high-pressure events determine recall. High pressure means $\omega\geq0.85$. We report planned and observed selections, recall, precision, and the share of measured paid occupancy captured. Paired intervals retain the three clock times together, resample seven consecutive target dates, and use 1-, 3-, and 5-km spatial cells.

\textit{Results.}
Table~\ref{tab:parking_validation} shows small magnitude-error reductions. At 500 m, identity Fixed-MTGNN changes December ratio MAE from 0.190803 to 0.190191, a difference of $-0.000612$ with a 95\% joint temporal and 3-km spatial interval of $[-0.001181,-0.000275]$. January changes from 0.191378 to 0.190885, with difference $-0.000494$ and interval $[-0.000995,-0.000180]$. The relative reductions are approximately 0.32\% and 0.26\%. Shared-LSTM and the graph inputs produce similar parking errors.

\begin{table}[!htbp]
\centering\small
\caption{Seven-day-ahead paid-parking validation with a 500-m POI catchment. Recall is the percentage of observed high-pressure events selected under a 10\% monitoring budget. Neural-input rows report five-seed means.}
\label{tab:parking_validation}
\begin{tabular}{lrrrr}
\toprule
 & \multicolumn{2}{c}{December} & \multicolumn{2}{c}{January} \\
Forecast input & Ratio MAE & Recall (\%) & Ratio MAE & Recall (\%) \\
\midrule
Parking calendar & 0.1908 & 51.49 & 0.1916 & 50.87 \\
Parking seven-day lag & 0.2062 & 44.82 & 0.2064 & 45.68 \\
Parking ridge & 0.1908 & 53.23 & 0.1914 & 53.01 \\
+ NLinear-ridge & 0.1904 & 53.30 & 0.1911 & 52.97 \\
+ Shared-LSTM & 0.1902 & 53.31 & 0.1909 & 53.00 \\
+ Matched-LSTM & 0.1902 & 53.33 & 0.1909 & 52.97 \\
+ GraphGRU, Identity & 0.1902 & 53.35 & 0.1909 & 52.97 \\
+ GraphGRU, NAICS-2 & 0.1902 & 53.36 & 0.1909 & 52.98 \\
+ Fixed-MTGNN, Identity & 0.1902 & 53.30 & 0.1909 & 52.96 \\
+ Fixed-MTGNN, NAICS-2 & 0.1902 & 53.32 & 0.1909 & 52.96 \\
\bottomrule
\end{tabular}
\end{table}


Monitoring differences are smaller than their uncertainty. At the 10\% budget, NAICS-2 GraphGRU changes high-pressure recall by $0.127$ percentage points in December and $-0.037$ in January relative to parking-only ridge. The corresponding joint intervals are $[-0.536,0.651]$ and $[-0.518,0.418]$ percentage points. Thus, the modest improvement in payment-based occupancy error yields similar monitoring coverage. The external outcome supports a narrow predictive contribution of POI forecasts in this sample; physical occupancy and realized inspection benefits remain separate quantities.
